% Admissible Non-Injective Transitions as the Primitive of Physical Description
% Cosmochrony companion paper M
% Provides the axiomatic foundation cited by Q1, Bell, ENI, and the white paper.

\documentclass[11pt,a4paper]{article}

\usepackage{amsmath,amssymb,amsthm}
\usepackage{enumerate}
\usepackage{booktabs}
\usepackage{array}
\usepackage{hyperref}
\usepackage{doi}

\newtheorem{theorem}{Theorem}[section]
\newtheorem{lemma}[theorem]{Lemma}
\newtheorem{proposition}[theorem]{Proposition}
\newtheorem{corollary}[theorem]{Corollary}
\newtheorem{definition}[theorem]{Definition}
\newtheorem{remark}[theorem]{Remark}

\title{Admissible Non-Injective Transitions\\
as the Primitive of Physical Description}

\author{J\'er\^ome Beau\\
\small Independent Researcher\\
\small \texttt{jerome.beau@cosmochrony.org}}

\date{2026}

\begin{document}

  \maketitle

  \begin{abstract}
    Quantum mechanics is the effective theory of physically real unresolved
    admissible structures and their resolution.
    This paper establishes the minimal axiomatic foundation from which this
    claim follows: physical structure is derived from a single primitive,
    the local structure of admissible transitions between observable states.
    Four axioms govern this structure: local projective admissibility (A1),
    structural non-injectivity (A2), admissibility as non-premature selection (A3),
    and projection locking (A4).
    From A1 and A2 alone, we derive irreversibility, the arrow of time, and the
    existence of structural fluctuations, without any additional postulate.
    From A3, we derive the proto-state as a physically real unresolved configuration
    shared between two successive states, and show that admissibility forces the
    retention of phase coherence.
    We further prove that A3 forbids any admissible factorisation of the fibre,
    which forces a non-trivial commutator between conjugate generators;
    combined with minimality (A1) and Born--Infeld parity, this identifies the
    symmetry group of the admissible fibre as $\mathrm{Heis}_3(\mathbb{Z}/q\mathbb{Z})$
    for a prime $q$, and the fibre itself as the Weil representation $V_\rho$ —
    closing the identification $F_n \simeq V_\rho$ as a theorem rather than a remark.
    From A4, we derive the discrete character of quantum transitions as a consequence
    of the interaction between continuous Born--Infeld saturation and the discrete
    shell structure of the observable space.
    This paper provides the axiomatic foundation cited by the companion
    papers~\cite{Q1,Bell,ENI} and serves as the formal basis for the derivations
    carried out therein.
  \end{abstract}

  \paragraph{Keywords.}
    Non-injective projection; Emergence; Projective structure; Admissibility; Born–Infeld constraint; Information loss;
    Projection entropy; Quantum emergence; Phase coherence; SU(2) structure; Gauge emergence; Fiber structure;
    Projective temporal ordering; Irreversibility.

    \clearpage
    \tableofcontents

    \clearpage

  \section{Introduction}
  \label{sec:intro}

  The Cosmochrony programme~\cite{Cosmochrony} aims to derive the structure of
  physical description from a minimal relational foundation, without postulating
  a background spacetime, a Hilbert space, or an external substrate.
  The only primitive is the local structure of admissible transitions between
  observable states.

  The central question driving the framework is:

  \begin{quote}
    Given an observable state $O_{n-1}$, what states are accessible from it?
  \end{quote}

  The answer to this question — encoded in the admissibility structure — is
  sufficient to derive irreversibility, temporal order, wave-like propagation,
  phase coherence, and the discrete character of quantum transitions.
  No substrate, no background, and no quantum postulate is required as input.

  \subsection*{Relation to companion papers}

    This paper provides the axiomatic foundation for a chain of companion results.
    The companion paper~\cite{Bell} establishes that Bell-type factorisation
    fails generically under non-injective projection (A2), without invoking
    nonlocal dynamics.
    The companion paper~\cite{ENI} proves that non-injectivity is a structural
    necessity of any genuinely emergent description.
    The companion paper~\cite{Q1} derives phase coherence, the singlet correlator,
    the Tsirelson bound, and the Born rule as structural consequences of
    admissibility — using Axioms A1--A4 as stated here, combined with the concrete
    spectral results of the O-series.
    The white paper~\cite{Cosmochrony} develops the full phenomenological programme;
    the present axiomatic reformulation is a foundational companion to it, not yet
    integrated into its main text.

  \subsection*{Organisation}

    Section~\ref{sec:axioms} states the four axioms.
    Section~\ref{sec:primitive} clarifies the role of $\Pi$ and the index $n$.
    Section~\ref{sec:ni-consequences} derives the structural consequences of
    non-injectivity: irreversibility, novelty, and temporal order.
    Section~\ref{sec:admissibility} derives the proto-state, phase, coherence,
    and the Heisenberg structure of the admissible fibre from the admissibility
    principle.
    Section~\ref{sec:locking} derives projection locking and the discrete character
    of quantum transitions.
    Section~\ref{sec:entanglement} treats entanglement and Bell non-applicability.
    Section~\ref{sec:summary} summarises what is derived from what, and identifies
    the open directions.

    \clearpage

  \section{The Four Axioms}
  \label{sec:axioms}

  We state the four axioms from which the framework proceeds.
  They are not independent postulates chosen for convenience — each captures
  a structural feature that is either necessary for genuine emergence~\cite{ENI}
  or forced by the requirement of bounded effective descriptions~\cite{BornInfeld}.

  \begin{definition}[Admissible transition]
    \label{def:admissible-transition}
    A transition from observable state $O_{n-1}$ to observable state $O_n$ is
    \emph{admissible} if it satisfies Axioms A1--A4 below.
  \end{definition}

  \begin{description}
    \item[\textbf{A1 — Local projective admissibility.}]
    For any observable state $O_{n-1}$, there exists a non-empty set $F_n$ of
    admissible successor directions: configurations of $O_{n-1}$ compatible with
    a transition to a stable next state.
    The set $F_n$ is defined by internal structural constraints
    (Born--Infeld bound, spectral envelope, symmetry requirements)
    and does not presuppose an ambient configuration space.

    \item[\textbf{A2 — Structural non-injectivity.}]
    Distinct admissible directions from $O_{n-1}$ may lead to the same resolved
    successor state $O_n$.
    The map $\Pi_n$ encoding this is generically non-injective.

    \item[\textbf{A3 — Admissibility as non-premature selection.}]
    An admissible transition preserves the full multiplicity of open directions
    until projection locking occurs.
    It does not select among them prematurely — that is, before the
    Born--Infeld saturation threshold is reached~\cite{BornInfeld}.

    \item[\textbf{A4 — Projection locking.}]
    Resolution occurs only when continuous Born--Infeld saturation meets the
    discrete shell support of the observable space.
    The transition from admissible to admitted is structurally discrete.
  \end{description}

  \begin{remark}[Independence of the axioms]
  A1 asserts existence of admissible directions.
  A2 asserts their generic non-injectivity.
  A3 constrains which transitions are admissible.
  A4 governs when resolution occurs.
  The four axioms are logically independent: none follows from the others.
  Their joint necessity is argued in Section~\ref{sec:primitive} and in the
  companion papers~\cite{ENI,BornInfeld}.
  \end{remark}

  \begin{remark}[Status of Born--Infeld in A3 and A4]
  The Born--Infeld saturation threshold appearing in A3 and A4 is not an
  additional postulate: it is the unique bounded completion of the effective
  action compatible with bounded projective flux, established in~\cite{BornInfeld}.
  A3 and A4 therefore do not introduce new structure beyond A1--A2 and the
  Born--Infeld companion.
  \end{remark}

  \begin{remark}[What projection locking selects]
    \label{rem:a4-selects}
    A4 resolves the transition to a unique observable state~$O_n$ in the image of~$\Pi_n$.
    It does not select a unique underlying configuration in~$F_n$.

    The multiplicity of admissible configurations in the fibre $F_n = \Pi_n^{-1}(O_n)$ is
    not eliminated by projection locking: distinct configurations $f, f' \in F_n$ with
    $\Pi_n(f) = \Pi_n(f') = O_n$ remain structurally present after locking, but become
    observationally inaccessible.
    The non-injectivity of~$\Pi_n$ therefore persists at every resolution step.

    This distinction --- observable uniqueness without configurational uniqueness --- prevents
    Bell-type factorisation~\cite{Bell} and forces the coherent, non-factorisable fibre
    structure derived in Section~\ref{subsec:heisenberg}.
  \end{remark}

  \section{The Primitive and the Role of $\Pi$}
  \label{sec:primitive}

  \subsection{The primitive}

    The only primitive of this framework is the pair $(O_{n-1}, F_n)$: an observable
    state and the set of admissible directions it admits.
    There is no external substrate, no hidden carrier, no background geometry.
    The fibre $F_n$ is not defined in an abstract ambient space — it is induced by
    $O_{n-1}$ itself through the admissibility constraints of A1--A3.

  \subsection{The projection $\Pi_n$ as descriptive abstraction}

    The notation $\Pi_n$ is a compact encoding of the local admissibility structure.
    It encodes three things simultaneously:
    \begin{enumerate}[(i)]
\item the admissible directions from $O_{n-1}$ — the fibre $F_n$ (A1, A2);
\item the Born--Infeld saturation bound — the limit on accessible amplitudes (A3);
\item projection locking — the constraint that resolution falls on an available
shell (A4).
    \end{enumerate}

    $\Pi_n$ is not a physical entity and not a dynamical map.
    It is a structural relation encoding the admissible identification between
    distinguishable configurations (elements of $F_n$) and indistinguishable
    observable outcomes (elements of $O_n$).
    Saying ``$\Pi_n$ is non-injective'' means that multiple directions from
    $O_{n-1}$ can lead to the same $O_n$ — a fact about the admissible transition
    structure, not a property of an external map.
    Non-injectivity does not originate from an underlying substrate, but from
    the structure of admissible configurations and their interaction-dependent
    realisation map.

  \subsection{The index $n$}

    The index $n$ does not represent a pre-existing time parameter.
    It counts transitions: the number of times a state has been succeeded by
    another under the admissibility constraints.
    Its temporal meaning is derived in Section~\ref{sec:ni-consequences} —
    it is not assumed.

  \subsection{Where non-injectivity lives}
    \label{subsec:ni-locus}

    A natural question arises: where does the non-injectivity of $\Pi_n$ come from,
    given that no underlying substrate is postulated?
    The answer is precise.

    Non-injectivity is a property of the interaction-dependent realisation map:
    the map from the space $F_n$ of structurally compatible configurations to the space
    $O_n$ of observable outcomes.
    When two distinct admissible configurations $f, f' \in F_n$ give rise to
    the same observable outcome $o \in O_n$, it is not because they are
    ``the same at a deeper level'' — it is because the interaction that
    realises the transition cannot distinguish them within $O_n$.
    The limit on distinguishability is set by the Born--Infeld capacity
    constraint (A3): the observable space $O_n$ has finite projective capacity,
    and configurations that differ only below this threshold are identified.

    Three properties of this realisation map are essential:
    \begin{enumerate}[(i)]
\item \emph{Interaction-dependence}: $\Pi_n$ encodes the interaction;
different interactions define different realisation maps and different
fibres $F_n$.
\item \emph{Internal definition of admissibles}: $F_n$ is defined by
internal constraints (BI bound, spectral envelope, symmetry) — not by
reference to an ambient substrate.
\item \emph{Minimal structure of $F_n$}: the admissibility constraints
force $F_n$ to carry a specific algebraic structure — SU(2) in the
spin-$\tfrac{1}{2}$ sector (O23) — which is derived from the constraints,
not assumed.
    \end{enumerate}

    The non-injectivity is therefore neither a property of a hidden substrate
    nor an arbitrary assumption: it is a structural consequence of finite
    projective capacity applied to an internally defined space of admissible
    configurations.

    \begin{remark}[The admissible fibre as equivalence class under projection]
      \label{rem:shadow-analogy}
      A useful image for the non-injectivity of~$\Pi_n$ is that of shadows cast by a fixed
      projection setup.
      A given shadow does not correspond to a single object:
      it corresponds to the set of all objects whose projection yields that shadow ---
      the admissible fibre $F_n = \Pi_n^{-1}(O_n)$.

      The directionality is essential.
      The fibre is not the collection of all shadows produced by a given object under varying
      projections (one configuration, many images).
      It is the collection of all configurations that produce the same observable under a fixed,
      non-injective~$\Pi_n$ (many configurations, one observable).
      This many-to-one structure is the source of irreversibility
      (Section~\ref{sec:ni-consequences}) and of the non-factorisability of the proto-state
      (Section~\ref{subsec:heisenberg}).
    \end{remark}

    \begin{remark}[Connection to the O-series: $F_n$ and $V_\rho$]
      \label{rem:fn-vrho}
      In the concrete spectral realisation of the O-series, the admissible fibre
      $F_n$ is realised as a spectral sector $V_\rho$ of the Heisenberg Cayley
      graph $\mathrm{Cay}(\mathrm{Heis}_3(\mathbb{Z}/q\mathbb{Z}), S_q)$.
      The sectors $V_\rho$ are not pre-existing configuration spaces: they are
      the structures induced by the admissibility constraints — the Born--Infeld
      bound, the spectral envelope, and the symmetry requirements of A1.
      In this sense, $V_\rho$ is the concrete realisation of $F_n$, not an
      ambient space in which $F_n$ is embedded:
      \begin{equation}
        F_n \;\simeq\; V_\rho.
      \end{equation}
      This identification, stated here as a structural connection to the O-series,
      is established as a theorem in Section~\ref{subsec:heisenberg}:
      the axioms A1--A3 and Born--Infeld parity together force the symmetry group
      of $F_n$ to be $\mathrm{Heis}_3(\mathbb{Z}/q\mathbb{Z})$, and the
      Stone--von Neumann theorem then identifies $F_n \simeq V_\rho$ uniquely.
    \end{remark}

    \begin{remark}[The two levels of $\Pi_n$: kernel and image (O24)]
      \label{rem:pi-two-levels}
      The non-injective projection $\Pi_n$ operates at two structurally distinct
      levels, established in O24~\cite{O24}:
      \begin{enumerate}[(i)]
  \item \emph{Kernel level} ($\ker \Pi_n$): governs microscopic multiplicity
  and fluctuation amplitude.
  Multiple configurations in $\Omega_n^{(c)}$ may project to the same
  observable outcome; their differences are invisible to admissible
  observables and appear as projection residuals.
  The dependence on $O_{n-1}$ — via the baseline $V_{n-1}^{(c)}$ — enters
  at this level: it determines which directions in $\Omega_n^{(c)}$ are
  new relative to the already realized subspace.
  \item \emph{Image level} ($\mathrm{Im}\,\Pi_n \cap \mathrm{Ntrl}$):
  governs the observable structure.
  Its real rank is rigidly fixed to 3 by the quaternionic maximality of
  $\mathrm{Im}\,\mathbb{H}$ (O23, Theorem~3.2) and is independent of the
  cardinality of $\ker \Pi_n$.
      \end{enumerate}
      The Verticality Lemma of O24 (Lemma~3.1) establishes that every
      Born--Infeld-admissible symmetry acts vertically: it maps each fibre
      $\Pi_n^{-1}(o)$ to itself and does not increase the rank of the image.
      The structural consequence is:
      \begin{quote}
        \emph{Non-injectivity affects degeneracy, not structure.}
      \end{quote}
      Enlarging $\ker \Pi_n$ — admitting more configurations in each fibre —
      increases the amplitude of quantum fluctuations but does not alter the
      observable rank, the cascade exponent $\delta_{\mathrm{pair}}$, or any
      other structural invariant of the admissible sector.
      The two levels of $\Pi_n$ are therefore decoupled: kernel structure is
      physically relevant for fluctuations; image structure is physically
      relevant for observables.
    \end{remark}


    \label{subsec:five-levels}

    The abstract structure of Section~\ref{sec:primitive} has a concrete realisation
    in the O-series that answers precisely where $\Omega_n$ lives and how it is
    constructed.
    The realisation proceeds through five levels, each induced by the previous one.

    \paragraph{Level 1 — The relational group $G_q$.}
      The combinatorial substrate of the BFS exploration is the discrete Heisenberg
      group
      \begin{equation}
        G_q \;=\; \mathrm{Heis}_3(\mathbb{Z}/q\mathbb{Z}),
        \qquad |G_q| = q^3,
      \end{equation}
      equipped with the standard symmetric generating set $S_q = \{X^{\pm 1}, Y^{\pm 1}\}$.
      $G_q$ is not a physical space: it is a relational structure, encoding how
      configurations relate to one another via the generators.
      The BFS explores $G_q$ in shells $S_n = B_n \setminus B_{n-1}$, where
      $B_n = \{g \in G_q : d(e,g) \leq n\}$.
      The shell structure $\{S_n\}$ is the concrete realisation of the
      transition-counting index $n$ of Section~\ref{sec:primitive}.

    \paragraph{Level 2 — The Weil representation $\rho_c$ and its ambient space.}
      For each non-zero character $c \in (\mathbb{Z}/q\mathbb{Z})^\times$, the
      discrete Heisenberg group $G_q$ admits an irreducible $q$-dimensional Weil
      representation~\cite{O12}:
      \begin{equation}
        \rho_c : G_q \;\longrightarrow\; \mathrm{GL}(\mathbb{C}^q),
        \qquad
        (\rho_c(a,b,\gamma)\,f)(x) \;=\;
        e^{2\pi i c(\gamma + bx)/q}\, f(x-a),
      \end{equation}
      acting on $L^2(\mathbb{Z}/q\mathbb{Z}) \cong \mathbb{C}^q$.
      The ambient space $\mathbb{C}^q$ is not a physical configuration space
      presupposed by the theory: it is the natural carrier of the irreducible
      representation of $G_q$ of central character $c$.
      It is induced by $G_q$ and $c$, not by any external assumption.

    \paragraph{Level 3 — The orbit span $\Omega_n$ as emerging subspace.}
      Given the initial vector $v_0 \in \mathbb{C}^q$ and the Weil action $\rho_c$,
      the orbit span at BFS depth $n$ is:
      \begin{equation}
        \label{eq:omega-n}
        \Omega_n^{(c)}
        \;:=\;
        \mathrm{span}\bigl\{\rho_c(g)\,v_0 : g \in B_n\bigr\}
        \;\subseteq\; \mathbb{C}^q.
      \end{equation}
      By construction, $\Omega_n^{(c)}$ is a minimal extension of the previously
      realized subspace:
      \begin{equation}
        \label{eq:omega-n-inclusion}
        V_{n-1}^{(c)} \;\subseteq\; \Omega_n^{(c)} \;\subseteq\; \mathbb{C}^q,
      \end{equation}
      where $V_{n-1}^{(c)} = \mathrm{GS}(B_{n-1})$ is the Gram--Schmidt span
      accumulated over the preceding shells.
      $\Omega_n^{(c)}$ is not defined by $O_{n-1}$, but is constrained by the
      already realized subspace $V_{n-1}^{(c)}$: the directions in
      $\Omega_n^{(c)} \setminus V_{n-1}^{(c)}$ are the candidates for new
      observable structure at depth $n$.

      \begin{definition}[Structurally compatible continuations]
        \label{def:omega-n}
        $\Omega_n^{(c)}$ is the space of \emph{structurally compatible continuations}
        from $O_{n-1}$: it contains every direction reachable by the relational
        structure of $G_q$ via $\rho_c$, without any admissibility filter applied.
        Admissibility is \emph{not} a property of $\Omega_n^{(c)}$;
        it is a property of $\Pi_n$ acting on $\Omega_n^{(c)}$.
      \end{definition}

      Three properties follow immediately:
      \begin{enumerate}[(i)]
\item \emph{Emergence}: $\Omega_n^{(c)}$ is not given a priori — it is built
incrementally as the BFS explores $G_q$.
\item \emph{Indirect conditioning on $O_{n-1}$}:
$\Omega_n^{(c)}$ is defined by the group action alone — the definition
\eqref{eq:omega-n} refers only to $B_n$ and $\rho_c$.
The dependence on $O_{n-1}$ is indirect: $O_{n-1}$ determines the
Gram--Schmidt span $\mathrm{GS}_{n-1}$ already accumulated over $B_{n-1}$,
which acts as the baseline against which new directions in $\Omega_n^{(c)}$
are measured.
This baseline is used by $\Pi_n$ (not by $\Omega_n^{(c)}$ itself) to
identify genuinely new directions, counted by $\sigma_c(n) = \delta r_n / |S_n|$~\cite{O12}.
In summary: $\Omega_n^{(c)}$ is determined by $(G_q, \rho_c, B_n)$;
the role of $O_{n-1}$ enters only through $\Pi_n$.
\item \emph{Pre-admissibility character}:
$\Omega_n^{(c)}$ contains all directions generated by the group action
of $G_q$ up to depth $n$ — including directions that $\Pi_n$ will
subsequently reject as inadmissible.
The admissibility filter is entirely in $\Pi_n$, not in $\Omega_n^{(c)}$.
      \end{enumerate}

    \paragraph{Level 4 — The admissibility filter $\Pi_n$.}
      The projection $\Pi_n$ acts on $\Omega_n^{(c)}$ via the BI admissibility
      constraint $A_n \leq c_\chi/\sqrt{\lambda_n}$ and the Gram--Schmidt
      orthogonalisation~\cite{O12,O22}.
      It selects, from all directions in $\Omega_n^{(c)}$, those compatible with
      the admissibility constraints.
      $\Pi_n$ encodes the laws revealed by the O-series — Born--Infeld bound,
      parity involution, projection locking — but does not define $\Omega_n^{(c)}$.

    \paragraph{Level 5 — The observable outcome $O_n$.}
      The observable state $O_n$ is the image of $\Pi_n$ applied to the admissible
      part of $\Omega_n^{(c)}$.
      It is a discrete, shell-indexed element of the observable
      space~\cite{O22}:
      \begin{equation}
        O \;=\; \bigsqcup_{n \in \mathbb{N}_q} O_n.
      \end{equation}

      \begin{remark}[Summary: the five-level hierarchy]
  The five levels constitute a single construction with no external inputs:
  \begin{equation}
    G_q
    \;\xrightarrow{\rho_c}\;
    \mathbb{C}^q
    \;\xrightarrow{\mathrm{BFS}}\;
    \Omega_n^{(c)}
    \;\xrightarrow{\Pi_n}\;
    O_n,
    \qquad
    V_{n-1}^{(c)} \subseteq \Omega_n^{(c)} \subseteq \mathbb{C}^q.
  \end{equation}
  $G_q$ is the relational structure.
  $\mathbb{C}^q$ is its representation carrier.
  $\Omega_n^{(c)}$ is the emergent space of structurally compatible continuations
  (not yet filtered by admissibility),
  constrained from below by the already realized subspace $V_{n-1}^{(c)}$.
  $\Pi_n$ is the admissibility filter.
  $O_n$ is the resolved observable state.
  No level presupposes the next: each is induced by the preceding structure
  and the admissibility constraints.
  This is the concrete answer to the question of where $\Omega_n$ lives:
  in $\mathbb{C}^q$, itself induced by $G_q$ — not in an external configuration
  space.
      \end{remark}

      \begin{remark}[The conjugate pair and $\Omega_n$ without preferred direction]
  The parity involution $c \leftrightarrow q-c$ (O18) acts on $\Omega_n^{(c)}$
  via complex conjugation~\cite{O17}:
  $\Omega_n^{(q-c)} = \overline{\Omega_n^{(c)}}$.
  The minimal admissible fibre is therefore the pair
  $\{V_n^{(c)}, V_n^{(q-c)}\} = \{V_n^{(c)}, \overline{V_n^{(c)}}\}$,
  with no preferred orientation — consistent with the admissibility principle
  of A3.
      \end{remark}

  \subsection{From $\mathbb{C}^q$ to continuous spacetime: the large-$q$ limit}
  \label{subsec:continuum-limit}

  The five-level hierarchy of Section~\ref{subsec:five-levels} is formulated
  for finite prime $q$.
  Physical spacetime is continuous.
  This subsection identifies the structural bridge between the two, without
  claiming to complete the derivation.

  \paragraph{The discrete structure at finite $q$.}
    At finite $q$, the observable space $O = \bigsqcup_{n \in \mathbb{N}_q} O_n$
    is discrete and finite.
    The ambient representation space $\mathbb{C}^q$ has dimension $q$.
    Physical quantities are functions on $\mathbb{Z}/q\mathbb{Z}$.

  \paragraph{The large-$q$ limit.}
    As $q \to \infty$ with the prime structure retained, three things happen
    simultaneously:
    \begin{enumerate}[(i)]
\item \emph{The representation space $\mathbb{C}^q$ approaches $L^2(\mathbb{R})$}.
Functions on $\mathbb{Z}/q\mathbb{Z}$ with fixed Fourier support become
functions on $\mathbb{R}$ in the Pontryagin dual limit.
The Weil representation $\rho_c$ of $G_q$ approaches the metaplectic
representation of the real Heisenberg group $\mathrm{Heis}_3(\mathbb{R})$.
\item \emph{The BFS shell structure approaches a continuous foliation}.
The polynomial ball growth $|B_n| \sim c\, n^4$ (Bass--Guivarc'h,
homogeneous dimension $D = 4$) induces, in the $q \to \infty$ limit,
a continuous layering of the configuration space compatible with
a four-dimensional effective geometry~\cite{Cosmochrony}.
\item \emph{The discrete admissibility filter $\Pi_n$ approaches a
differential constraint}.
The Born--Infeld admissibility bound $A_n \leq c_\chi/\sqrt{\lambda_n}$
becomes, in the continuum limit, the Born--Infeld Lagrangian constraint
on effective field configurations~\cite{BornInfeld}.
    \end{enumerate}

    The continuum spacetime observable is therefore not a separate postulate:
    it is the large-$q$ limit of the discrete observable space $O$, with the
    four-dimensional structure inherited from the homogeneous dimension of $G_q$.

    \begin{remark}[Status of the continuum limit]
      \label{rem:continuum-limit}
      The convergence $\mathbb{C}^q \to L^2(\mathbb{R})$ and the emergence of
      continuous spacetime geometry from the BFS shell structure are established
      qualitatively in~\cite{Cosmochrony}.
      The companion paper~\cite{Q5a} initiates the rigorous treatment of this
      limit, formulating it in the category of Hilbert spaces and operator algebras
      and decomposing it into three distinct convergences: a Hilbert inductive limit
      (T1), a representational convergence of Weil generators to metaplectic
      generators (T2), and a Mosco convergence of admissibility forms to a
      second-order operator $L_\Pi$ on $L^2(\mathbb{R})$ (T3).
      The continuum limit is reduced to a single spectral tightness condition
      (Conjecture~C of~\cite{Q5a}), which is confirmed quantitatively across all
      tested primes: the spectral tail mass of admissible test vectors satisfies
      $E_q(q/3)/\|f_q\|^2 < 2.1 \times 10^{-5}$ uniformly, with values decreasing
      toward machine precision as $q$ grows.
      The embedding hypothesis~[H1] is proved in the low-frequency regime relevant
      to admissible sequences; the Weil-block compatibility condition and the
      identification of the effective metric from the symbol of $L_\Pi$ are the
      subject of the companion paper Q5b.
      The present remark identifies the mechanism; \cite{Q5a} establishes the
      operator-theoretic framework and the numerical evidence.
    \end{remark}

    \begin{remark}[Why four dimensions]
  The homogeneous dimension $D = 4$ of the Heisenberg group $G_q$
  (Bass--Guivarc'h theorem: $|B_n| \sim n^4$) is not chosen to match
  the observed four dimensions of spacetime.
  It is a structural consequence of the nilpotency class and rank of
  $\mathrm{Heis}_3$: the group has three generators with one non-trivial
  commutator, giving $D = 2 \cdot \mathrm{rank} + \mathrm{class} = 4$.
  The coincidence with the observed dimension of spacetime is therefore
  a structural prediction of the framework, not a parameter fit.
    \end{remark}

    \clearpage

  \section{Structural Consequences of Non-Injectivity}
  \label{sec:ni-consequences}

  Non-injectivity (A2) and the existence of distinct successor states together
  imply three structural facts whose sources are distinct.

  \subsection{Irreversibility}

    \begin{proposition}[Irreversibility from non-injectivity]
      \label{prop:irreversibility}
      Under A2, every admissible transition involves an irreversible loss of
      structural information.
    \end{proposition}

    \begin{proof}
      Since multiple directions from $O_{n-1}$ map to the same $O_n$ under
      $\Pi_n$, the information distinguishing these directions is absent from
      $O_n$.
      No subsequent transition $\Pi_{n'}$ for $n' > n$ can recover this
      information: the map $\Pi_n$ has positive projection entropy
      $S_{\Pi_n} > 0$~\cite{ENI}, and no admissible transition can reduce
      $S_\Pi$ below zero.
      The loss is therefore structural and permanent.
    \end{proof}

  \subsection{Novelty and structural fluctuations}

    The existence of a distinct successor state $O_n \neq O_{n-1}$ defines a
    structural residue:
    \begin{equation}
      \label{eq:delta-n}
      \Delta_n \;:=\; \mathrm{Res}(O_{n-1}, O_n;\, \Pi_n),
    \end{equation}
    the measure of how much $O_n$ is irreducible to a continuation of $O_{n-1}$
    under $\Pi_n$.
    That $\Delta_n$ is non-trivial is not an additional assumption — it is the
    condition under which we index $O_n$ as a distinct state.
    These residues are what physical descriptions call fluctuations or quantum noise.

  \subsection{Temporal order}

    \begin{theorem}[Temporal order from A2]
      \label{thm:temporal-order}
      Under A2, the chain of observable states $\{O_n\}_{n \geq 0}$ carries a
      canonical oriented, non-trivial partial order.
    \end{theorem}

    \begin{proof}
      Irreversibility (Proposition~\ref{prop:irreversibility}) establishes an
      asymmetric relation: $O_{n-1} \to O_n$ is not reversible, so $O_{n-1}$
      and $O_n$ are not equivalent under the transition structure.
      Novelty ($\Delta_n \neq 0$) ensures $O_n \neq O_{n-1}$, so the order
      is non-trivial.
      The composition of transitions defines transitivity.
      Together: asymmetry, irreflexivity, and transitivity give a strict partial
      order compatible with the chain structure.
    \end{proof}

    \begin{remark}[Time as derived structure]
  The arrow of time is not postulated.
  It is the oriented order generated by the chain of non-injective transitions
  under A2.
  The index $n$ acquires its temporal meaning from this order — not the reverse.
  The quantitative structure of this projective temporal ordering — cumulative
  projected capacity and its numerical validation from the spectral admissibility
  cascade — is developed in~\cite{Beau2026pto}.
    \end{remark}

  \section{Admissibility, Wave, and Phase Coherence}
  \label{sec:admissibility}

  \subsection{The admissibility principle}

    Axiom A3 introduces a physical constraint on transitions:

    \begin{quote}
      A transition is admissible if and only if it preserves the full multiplicity
      of open admissible directions without selecting among them prematurely.
    \end{quote}

    Premature selection means making a determinative choice that nothing in the
    current state $O_{n-1}$ justifies.
    Admissibility forbids it: the transition must remain structurally neutral among
    all open directions until projection locking (A4) forces resolution.

  \subsection{The proto-state}
    \label{subsec:proto-state}

    \begin{definition}[Proto-state]
      \label{def:proto-state}
      The \emph{proto-state} $\Psi_n$ associated with a transition from $O_{n-1}$
      to $O_n$ is the unresolved configuration shared between $O_{n-1}$ and $O_n$:
      the carrier of the transition, belonging fully to neither endpoint.
    \end{definition}

    The proto-state is physically real — it is the configuration of the system
    during the transition, before projection locking.
    It is not an epistemic representation of uncertainty about a definite state:
    it is the actual structural state of the system between two realized
    observable states.

    \begin{proposition}[Structural indeterminacy of the proto-state]
      \label{prop:indeterminacy}
      Under A1 and A3, the proto-state $\Psi_n$ is genuinely indeterminate among
      the admissible directions in $F_n$.
    \end{proposition}

    \begin{proof}
      By A1, $F_n$ is non-empty.
      By A3, the transition does not select prematurely among the elements of $F_n$.
      Therefore $\Psi_n$ is not identified with any particular element of $F_n$
      before projection locking: it is open across all of $F_n$ simultaneously.
      This is structural indeterminacy, not epistemic ignorance.
    \end{proof}

  \subsection{Wave, phase, and coherence}

    The proto-state, being structurally open across all elements of $F_n$
    simultaneously, propagates through all admissible directions at once.
    This is the physical origin of wave-like behaviour: not a metaphor, but a
    direct consequence of admissibility.

    A configuration propagating through multiple directions simultaneously carries
    relative phase structure between those directions.
    Phase is therefore not postulated — it is a consequence of structural
    indeterminacy.

    Admissibility (A3) preserves this phase structure: premature selection would
    collapse the relative phases, but A3 forbids it.
    The chain of consequences is therefore:

    \begin{equation}
      \begin{aligned}
        \text{A3}
        &\;\Longrightarrow\; \text{structural indeterminacy} \\
        &\;\Longrightarrow\; \text{proto-state (wave)} \\
        &\;\Longrightarrow\; \text{phase} \\
        &\;\Longrightarrow\; \text{coherence preserved.}
      \end{aligned}
    \end{equation}

    \begin{remark}[The admissible fibre as precursor of Hilbert structure]
  The fibre $F_n$, equipped with the phase structure forced by A3, is the
  natural precursor of an effective Hilbert sector: it carries complex
  amplitudes and interference structure.
  This is not yet a full Hilbert space — the inner product, completeness,
  and probability measure require additional derivation (see Q1, Q2 in
  Section~\ref{sec:summary}).
  But the amplitude structure and interference follow from A3 alone.
    \end{remark}

    \begin{remark}[Double-slit experiment]
  In the double-slit experiment, the proto-state propagates through both slits
  because the admissibility constraints do not select between them:
  both are open directions in $F_n$.
  The interference pattern is a direct consequence of coherent propagation
  through all admissible directions.
  The impact on the screen is the moment the proto-state reaches the
  Born--Infeld saturation threshold — projection locking (A4) forces it onto
  an available shell, resolving it into a localized state $O_n$.
    \end{remark}

  \subsection{Non-factorisability and the emergence of Heisenberg structure}
    \label{subsec:heisenberg}

    Axiom A3 does not only enforce phase coherence.
    It also constrains the algebraic structure of the admissible fibre $F_n$ itself,
    forcing it to carry a Heisenberg structure rather than an arbitrary unitary
    symmetry.

    \begin{remark}[The structural tension between A2 and A4]
      \label{rem:a2-a4-tension}
      The non-factorisability established in this section can be understood as the minimal
      consistent resolution of a structural tension between A2 and A4.

      A2 asserts that the multiplicity of admissible configurations in~$F_n$ cannot be
      eliminated: distinct configurations lead to the same observable, and this
      indistinguishability is structural, not epistemic.
      A4 asserts that the observable outcome is unique at locking.
      Together, they impose an irresolvable constraint: the multiplicity cannot disappear
      (A2 would be violated), yet it cannot remain directly observable (A4 would be violated).

      A3 --- admissibility as non-premature selection --- is the operative axiom that
      determines how this tension resolves.
      It forbids any factorisation of~$F_n$ that would make the multiplicity distinguishable
      prior to locking, forcing it to organise into a globally coherent, irreducible structure.
      Lemma~\ref{lem:non-factorisation} formalises this constraint;
      Theorem~\ref{thm:heisenberg} identifies the resulting structure as Heisenberg type.
      Quantum structure is therefore the minimal algebraic form compatible with the joint
      constraint of A2, A3, and A4.
    \end{remark}

    \begin{lemma}[Admissible non-factorisability]
      \label{lem:non-factorisation}
      Under A1--A3, any decomposition
      \[
        F_n \;=\; F^{(1)} \oplus F^{(2)}
      \]
      that is stable under the admissible generators of $F_n$ is forbidden prior
      to projection locking.
      In particular, A3 forces the action of admissible generators on $F_n$ to
      be irreducible.
    \end{lemma}

    \begin{proof}
      A3 forbids premature selection among open directions.
      A decomposition of $F_n$ stable under admissible generators separates
      directions by invariants that are internal to the fibre — invariants carried
      by the generators themselves, which are constitutive of the admissibility
      structure by A1.
      Such a decomposition therefore resolves the proto-state partially, by
      assigning each open direction to a definite sector, prior to any
      Born--Infeld saturation event.
      This contradicts A3.
      Irreducibility follows immediately: any proper stable sub-representation
      would yield such a forbidden decomposition.
    \end{proof}

    \begin{theorem}[A1--A3 and BI parity force Heisenberg structure]
      \label{thm:heisenberg}
      Under A1--A3 and Born--Infeld parity, the symmetry group of the admissible
      fibre $F_n$ is isomorphic to $\mathrm{Heis}_3(\mathbb{Z}/q\mathbb{Z})$ for
      some prime $q$.
      Consequently, $F_n$ is isomorphic to the Weil representation $V_\rho$ on
      $L^2(\mathbb{Z}/q\mathbb{Z})$.
    \end{theorem}

    \begin{proof}[Proof sketch]
      Let $X$ and $\sigma(X)$ be the conjugate generators of $F_n$ under
      Born--Infeld parity.
      If $[X, \sigma(X)] = 0$, then $F_n$ decomposes into simultaneous eigenspaces,
      yielding a stable decomposition forbidden by Lemma~\ref{lem:non-factorisation}.
      Hence $[X, \sigma(X)] \neq 0$.
      Minimality (A1) forces the commutator $Z := [X, \sigma(X)]$ to be central:
      any non-central $Z$ would extend the algebra beyond the conjugate pair,
      contradicting the minimal generative closure forced by A1.
      The group $\langle X, \sigma(X) \rangle$ is then a non-trivial central
      extension of $(\mathbb{Z}/q\mathbb{Z})^2$ by $\langle Z \rangle$, isomorphic
      to $\mathrm{Heis}_3(\mathbb{Z}/q\mathbb{Z})$ by the classification of such
      extensions.
      Lemma~\ref{lem:non-factorisation} further implies that $|Z|$ is prime: if
      $|Z|$ were composite, the representation of $F_n$ would decompose over the
      characters of $\langle Z \rangle$, yielding a forbidden stable decomposition.
      Finally, Born--Infeld parity forces $c \neq 0$~\cite{BornInfeld}, so the
      central character is non-trivial; the discrete Stone--von Neumann theorem then
      identifies $F_n \simeq V_\rho$ uniquely.
      The detailed proof of each step is given in the companion
      paper~\cite{HeisenbergStructure}.
    \end{proof}

    \begin{equation}
      \begin{aligned}
        \text{A3}
        &\;\Longrightarrow\; \text{no admissible factorisation}
        \;\Longrightarrow\; [X, \sigma(X)] \neq 0
        \;\Longrightarrow\; \text{central extension,} \\
        \text{A1}
        &\;\Longrightarrow\; \text{minimal generative closure}
        \;\Longrightarrow\; \text{Heisenberg algebra,} \\
        \text{A3}
        &\;\Longrightarrow\; \text{irreducibility}
        \;\Longrightarrow\; |Z| = q \text{ prime}
        \;\Longrightarrow\; \mathrm{Heis}_3(\mathbb{Z}/q\mathbb{Z}), \\
        \text{BI parity}
        &\;\Longrightarrow\; c \neq 0
        \;\Longrightarrow\; F_n \;\simeq\; V_\rho.
      \end{aligned}
    \end{equation}

    \begin{remark}[Uncertainty principle as a corollary]
      \label{rem:uncertainty}
      Theorem~\ref{thm:heisenberg} subsumes the open direction~Q6 of
      Section~\ref{sec:summary}.
      The incompatibility of conjugate observables $X$ and $\sigma(X)$ — i.e.\
      the Heisenberg uncertainty principle — is a direct corollary: it follows
      from $[X, \sigma(X)] \neq 0$ together with the standard operator inequality.
      The uncertainty principle is therefore not a property of the observables,
      but a property of the non-resolved proto-state: it is the algebraic
      signature of the prohibition on premature selection encoded in A3.
    \end{remark}

    \begin{remark}[Canonical quantisation as structural necessity]
      The relation $[X, \sigma(X)] = Z$ defines a minimal discrete symplectic
      structure, of which the Heisenberg algebra is the canonical quantisation.
      Theorem~\ref{thm:heisenberg} therefore establishes that canonical
      quantisation is not a modelling choice: it is the unique algebraic structure
      compatible with the requirement that the proto-state remain genuinely
      non-resolved until projection locking.
    \end{remark}



    A quantum superposition is an unresolved proto-state: a physically real
    configuration whose fibre maintains coherent relations between multiple open
    admissible directions.
    There is no ontological paradox — superposition is the normal condition of
    a system between two transitions.

    A measurement is a transition whose admissibility constraints are strong enough
    to force resolution of the proto-state onto a single admissible direction.
    Free evolution and measurement differ in the strength of their admissibility
    constraints, not in their fundamental mechanism.

    In this sense, quantum mechanics is the effective theory of physically real
    unresolved admissible structures and their resolution.

  \subsection{The wavefunction as effective representation of the proto-state}
    \label{subsec:wavefunction}

    A quantum wave is not a physical substance propagating in space.
    It is the effective representation of a physically real unresolved proto-state.
    It encodes the coherent structure of admissible directions prior to projection
    locking.

    More precisely: the proto-state $\Psi_n$ is real;
    the wavefunction $\psi$ is a coordinate description of the coherent admissible
    structure carried by $\Psi_n$ in the spectral realisation $V_\rho$ of $F_n$.
    When we write $\psi$, we are not describing a substance propagating in space
    — we are encoding the relational phase structure of an unresolved
    configuration in the space of admissible directions.
    Quantum mechanics does not describe realized physical states:
    it describes states in the process of being resolved.

    This resolves the standard interpretational difficulties without additional
    postulate:
    \begin{itemize}
      \item \emph{Superposition} is the normal condition of a proto-state — the
      simultaneous structural openness across all elements of $F_n$.
      \item \emph{Interference} is the direct consequence of coherent propagation
      through multiple admissible directions before resolution.
      \item \emph{Collapse} is not a physical process applied to a pre-existing state
      but the resolution of the proto-state at projection locking — the transition
      from unresolved admissible structure to a realized observable outcome.
    \end{itemize}

  \subsection{Why complex amplitudes}
    \label{subsec:complex-amplitudes}

    Given that the proto-state is a physically real, unresolved, coherent
    configuration, what is the minimal mathematical structure capable of
    representing it?

    Real positive weights are insufficient.
    They describe ignorance about which of several pre-existing alternatives is
    realized — a classical mixture.
    The proto-state is not a mixture: it is a single coherent configuration
    with definite phase relations between admissible directions.

    Real signed weights are also insufficient.
    They allow opposition of sign but do not support a rich enough phase structure
    under composition.
    In particular, they cannot represent the full coherence of a proto-state whose
    directions combine under the SU(2) structure forced by O23.

    The minimal structure that can:
    (i) represent simultaneous coherent openness across multiple directions,
    (ii) encode their relative phase relations,
    (iii) compose correctly under the SU(2) symmetry of the admissible fibre,
    (iv) support interference — the physical consequence of coherence before
    resolution —
    is a complex amplitude.

    Therefore: complex amplitudes are not introduced by convention.
    They are the minimal representation of a coherent unresolved proto-state.
    The imaginary unit encodes the relative phase between admissible directions;
    it enters because the proto-state is not a mixture but a coherent
    superposition.

  \subsection{Why the Born weight is the squared norm}
    \label{subsec:born-weight}

    Upon projection locking (A4), the unresolved proto-state is resolved onto
    a single admissible direction.
    The quantity associated with each direction at resolution — the
    \emph{intensity of realisation} — must satisfy four structural requirements:

    \begin{enumerate}[(i)]
\item \emph{Positivity}: a weight on admissible directions cannot be negative.
\item \emph{Global phase invariance}: multiplying $\psi$ by a global phase
$e^{i\theta}$ does not change the physical proto-state, so the weight must
be invariant under this transformation.
\item \emph{Normalisability}: the total weight over all admissible directions
must be finite and summable to one.
\item \emph{Coherence preservation before resolution}: the weight must not
erase phase relations prematurely — interference must remain possible up to
the moment of resolution.
    \end{enumerate}

    The squared norm $|\psi|^2$ is the unique minimal quantity satisfying all
    four requirements.
    The norm $|\psi|$ fails (ii) under phase shifts and (iv) — it introduces
    real-valued interference which is not the correct structure.
    Any higher power $|\psi|^p$ for $p \neq 2$ fails (iii) under composition
    of superpositions.
    Any non-analytic function of $\psi$ fails (iv).

    The Born weight is therefore not a postulate.
    It is the positive intensity associated with the resolution of a coherent
    admissible structure — the unique minimal such intensity compatible with
    the four structural requirements.

    \begin{remark}[Relation to the Born rule derivation in the companion paper]
  The argument above establishes the structural necessity of the squared norm
  as the resolution intensity of the proto-state.
  The companion paper~\cite{Q1} makes this precise as a theorem in the SU(2)
  parity sector: the Born rule $P(a=+1|\hat{a}) = |\langle +\hat{a}|\psi\rangle|^2$
  is the unique probability assignment compatible with the coherent fibre
  structure, the SU(2) observable algebra, and the derived singlet correlator.
  The two derivations are complementary: the present argument gives the
  structural motivation; the companion gives the uniqueness proof.
    \end{remark}

    \begin{remark}[The wavefunction lives in the admissible fibre, not in space]
  A common source of confusion is to picture the wavefunction as an object
  propagating in physical space.
  In the present framework, $\psi$ encodes the structure of $F_n$ — the
  admissible directions from $O_{n-1}$.
  Its concrete realisation is the spectral sector $V_\rho$ induced by the
  admissibility constraints on the Heisenberg Cayley graph (O-series).
  $F_n$ and $V_\rho$ are not pre-existing spaces: they are the spaces
  defined by the admissibility constraints.
  The wavefunction does not propagate in space; it encodes the relational
  phase structure of an unresolved admissible configuration.
    \end{remark}


    \label{sec:locking}

    Axiom A4 establishes that resolution is discrete.
    This discreteness is not postulated: it follows from the interaction of two
    structures already present in the framework.

  \subsection{Two types of discreteness}

    The Born--Infeld saturation condition~\cite{BornInfeld} operates on a
    continuous amplitude function.
    Its saturation locus $\mathcal{L}_{\mathrm{BI}} \subset \mathbb{R}$ is
    generically real-valued — its elements carry no information about shell
    structure.

    The observable space decomposes discretely:
    \begin{equation}
      O \;=\; \bigsqcup_{n \in \mathbb{N}_q} O_n,
    \end{equation}
    indexed by BFS shells of the Heisenberg Cayley graph~\cite{O22}.
    No component $O_{n'}$ exists for non-admissible depths.

    The proto-state is resolved at the first element of
    $\mathcal{L}_{\mathrm{BI}} \cap \mathbb{N}_q$ — the first depth where
    Born--Infeld saturation coincides with an available shell.
    This is the Projection Locking theorem~\cite{O22}.

    \begin{theorem}[Discreteness from projection locking, after O22]
      \label{thm:discrete-locking}
      Under A3 and A4, quantum transitions are discrete.
      The discreteness is a consequence of the interaction of continuous
      Born--Infeld saturation with the discrete admissibility filter of $\Pi_n$,
      not a postulate.
    \end{theorem}

  \subsection{Three stable directions}

    The number of independent stable admissible directions in the fibre is
    exactly three, derived algebraically from the admissibility constraints
    alone~\cite{O23}:
    \begin{equation}
      \begin{aligned}
        \text{BI parity}
        &\;\Rightarrow\; \text{fibre involution } c \leftrightarrow q-c \\
        &\;\Rightarrow\; \text{non-abelian admissible support} \\
        &\;\Rightarrow\; \mathbb{H}\text{-minimality} \\
        &\;\Rightarrow\; \dim_{\mathbb{R}} \mathrm{Im}\,\mathbb{H} = 3 \\
        &\;\Rightarrow\; \Sigma_c(n_3) = 3.
      \end{aligned}
    \end{equation}
    The emergence of exactly three stable directions — and with it, three
    particle generations — is algebraic in origin, not geometric or
    phenomenological.

  \section{Entanglement and Bell Non-Applicability}
  \label{sec:entanglement}

  \subsection{Entanglement as shared ancestral fibre}

    Two observable systems $A$ and $B$ are entangled when they share a common
    ancestral fibre: a past state $O_{n-k}$ in which they formed a single
    non-factorisable configuration.

    The shared fibre $F_{n-k}$ reduces the admissible directions for each
    subsystem: the admissible directions for $A$ are constrained by $B$,
    and vice versa.
    The correlations arise not from any interaction at the moment of
    measurement, but from the persistence of the shared fibre — a residue of
    the past transition in which $A$ and $B$ were joined.

  \subsection{Why Bell's factorisation assumptions do not apply}

    Bell's theorem~\cite{Bell1964} constrains models in which measurement outcomes
    are determined by a variable $\lambda$ satisfying:
    (i) $\lambda$ is local (assignable separately to $A$ and $B$),
    (ii) $\lambda$ is present (defined at the moment of measurement),
    (iii) $\lambda$ is independent (statistically independent of measurement settings).

    The fibre variable in the present framework satisfies none of these conditions.
    It is not local (it belongs to the shared fibre, not to a spatial region),
    not present (it belongs to a past transition, not the current state),
    and not independent ($A$ and $B$ share the same ancestral fibre,
    which cannot be factorised).

    The Bell factorisation assumption
    $P(a,b|\lambda) = P(a|\lambda) \cdot P(b|\lambda)$
    fails because it requires $\lambda$ to be assignable separately to $A$ and $B$:
    the shared fibre is structurally non-factorisable.
    Bell-type inequalities are therefore not violated in this framework —
    their preconditions do not hold.
    This is established formally in the companion paper~\cite{Bell}.

  \section{Summary: What Is Derived from What}
  \label{sec:summary}

  Table~\ref{tab:derived} summarises the derivation chain.

  \begin{table}[h]
    \centering
    \small
    \begin{tabular}{lll}
      \toprule
      \textbf{Concept} & \textbf{Source} & \textbf{Status}
      \\
      \midrule
      Irreversibility & Non-injectivity (A2) & Proved (Prop.~
      \ref{prop:irreversibility}) \\
      Structural fluctuations & Relational residue $\Delta_n$ & Derived
      \\
      Temporal order & Irreversibility + novelty (A2) & Proved (Thm.~
      \ref{thm:temporal-order}) \\
      Proto-state (wave) & Unresolved multiplicity (A1, A3) & Derived (Prop.~
      \ref{prop:indeterminacy}) \\
      Structural indeterminacy & Multiple open directions (A1, A3) & Derived
      \\
      Phase coherence & Admissibility preserves open directions (A3) & Derived
      \\
      Coherent amplitude structure & Admissible fibre + phase & Derived
      \\
      Interference & Coherent amplitude structure & Derived
      \\
      Discrete transitions & Projection locking (A4, O22) & Proved (Thm.~
      \ref{thm:discrete-locking}) \\
      Three stable directions & $\dim_\mathbb{R} \mathrm{Im}\,\mathbb{H} = 3$ (O23) & Proved
      \\
      Superposition & Unresolved proto-state & Derived
      \\
      Collapse & Proto-state resolution at locking & Derived
      \\
      Entanglement & Shared non-factorisable ancestral fibre & Derived
      \\
      Bell non-applicability & Temporal, non-local, non-independent fibre & Proved~\cite{Bell}
      \\
      Phase coherence $\Rightarrow$ QM & Admissibility + O-series & Proved~\cite{Q1}
      \\
      Born rule (SU(2) sector) & Correlator + coherence (Q1) & Proved~\cite{Q1}
      \\
      $F_n \simeq V_\rho$ (Heisenberg structure) & Non-factorisability (A1, A3) + BI parity &
      Proved (Thm.~\ref{thm:heisenberg})
      \\
      Uncertainty principle & $[X,\sigma(X)]\neq 0$ (Thm.~\ref{thm:heisenberg}) &
      Proved (Rem.~\ref{rem:uncertainty})
      \\
      \midrule
      Full Hilbert structure & Beyond SU(2) parity sector & \textbf{Open}
      \\
      Born rule (general) & Beyond SU(2) parity sector & \textbf{Open}
      \\
      \bottomrule
    \end{tabular}
    \caption{Derivation chain from the four axioms.
    ``Derived'' denotes a consequence that follows from the axioms by
    structural argument; ``Proved'' denotes a result with an explicit proof
    either in this paper or in a cited companion.}
    \label{tab:derived}
  \end{table}

  \subsection*{Open directions}

    \paragraph{Q2 — Born rule beyond SU(2).}
      The Born rule $P(a=+1|\hat{a}) = |\langle +\hat{a}|\psi\rangle|^2$ is derived
      in~\cite{Q1} for the SU(2) parity sector.
      A first step beyond this sector is taken in~\cite{Q2}, which derives the
      Born rule, singlet correlator, and Tsirelson bound for the co-admissible
      spin-$\tfrac{3}{2}$ sector on the binary icosahedral group $2I$, and
      identifies SU(2) as the unique stable sector under admissibility refinement
      in the large-$p$ limit.
      The full extension to arbitrary observables, multipartite systems, and
      continuous spectra requires extending the fibre structure beyond the
      two-element parity case and remains the primary open direction.

    \paragraph{Q5 — Continuum limit of the fibre.}
      Theorem~\ref{thm:heisenberg} establishes $F_n \simeq V_\rho$ at finite
      prime~$q$.
      The relationship between this discrete structure and continuous spacetime
      geometry in the large-$q$ limit — i.e.\ how $V_\rho$ on
      $L^2(\mathbb{Z}/q\mathbb{Z})$ approaches $L^2(\mathbb{R})$ and how the
      BFS shell structure produces a four-dimensional effective geometry —
      is addressed in the companion spectral geometry paper~\cite{Q5a}
      (Section~\ref{subsec:continuum-limit}).
      That paper reduces Q5 to a single spectral tightness condition (Mosco
      tightness of the admissibility forms), confirms this condition
      quantitatively across all tested primes, and proves the embedding in the
      low-frequency regime relevant to admissible sequences.
      The identification of the effective metric from the principal symbol of the
      limit operator $L_\Pi$ and the Lorentzian signature selection are the
      subject of Q5b.

  \section*{Acknowledgements}

  The author acknowledges the use of large language models as a supportive tool
  for refining language, structure, and internal consistency during the development
  of this manuscript.
  All conceptual contributions, theoretical choices, and interpretations remain
  the sole responsibility of the author.

  \begin{thebibliography}{99}

    \bibitem{HeisenbergStructure}
    J.~Beau,
    ``Non-Factorisability and the Emergence of Heisenberg Structure
    from Admissibility Constraints,''
    Cosmochrony companion paper, Zenodo (2026).
    \doi{10.5281/zenodo.19635395}.

    \bibitem{Cosmochrony}
    J.~Beau,
    ``Cosmochrony: A Pre-Geometric Relational Framework for Emergent Physics,''
    white paper, Zenodo (2026).
    \doi{10.5281/zenodo.17957509}.

    \bibitem{ENI}
    J.~Beau,
    ``Non-Injectivity as a Structural Necessity of Genuine Emergence:
    A No-Go Theorem for Faithful Emergent Descriptions,''
    Cosmochrony companion paper~L, Zenodo (2026).
    \doi{10.5281/zenodo.19391746}.

    \bibitem{BornInfeld}
    J.~Beau,
    ``Bounded Relaxation and the Dynamical Selection of Spacetime Geometry,''
    Cosmochrony companion paper~C, Zenodo (2026).
    \doi{10.5281/zenodo.18407505}.

    \bibitem{O12}
    J.~Beau,
    ``Exact Weil-Block Projective Capacity on Heisenberg Graphs: Resolving the Final Obstruction to $\delta$
    Extraction,''
    spectr. adm. subprg. O12 paper, Zenodo (2026).
    \doi{10.5281/zenodo.19194798}.

    \bibitem{O17}
    J.~Beau,
    ``Fibre-Level Admissibility from Conjugate Weil Blocks: Structural Derivation of the Pair Observable,''
    O16 spectr. adm. subprg. O17 paper, Zenodo (2026).
    \doi{10.5281/zenodo.19299611}

    \bibitem{Bell}
    J.~Beau,
    ``Non-Injective Projections and the Structural Failure of Bell Factorization,''
    Cosmochrony companion paper~B, Zenodo (2026).
    \doi{10.5281/zenodo.18371173}.

    \bibitem{Q1}
    J.~Beau,
    ``From Admissibility to Quantum Structure: Phase Coherence and Correlations from Non-Injective Projection,''
    Cosmochrony companion paper~Q1, Zenodo (2026).
    \doi{10.5281/zenodo.19561060}.

    \bibitem{Q2}
    J.~Beau,
    ``Quantum Structure beyond Spin-$\tfrac{1}{2}$: Co-Admissible Sectors and
    the Emergence of SU(2) as a Fixed Point,''
    Cosmochrony companion paper~Q2, Zenodo (2026).
    \doi{10.5281/zenodo.19616444}.

    \bibitem{O22}
    J.~Beau,
    ``Shell-Alignment from Projection Locking: A Discrete Admissibility Theorem
    under Born--Infeld Saturation,''
    O22 spectral admissibility sub-programme paper, Zenodo (2026).
    \doi{10.5281/zenodo.19367375}.

    \bibitem{O24}
    J.~Beau,
    ``Vertical Non-Injectivity and the Stability of the Observable Rank:
    Closing the Unconditional Transfer $c_\chi \to \delta_{\mathrm{pair}} \to \beta^*$,''
    O24 spectral admissibility sub-programme paper, Zenodo (2026).
    \doi{10.5281/zenodo.19386581}.

    \bibitem{O23}
    J.~Beau,
    ``Three Stable Directions from Quaternionic Minimality: Derivation of
    $\Sigma_c(n_3) = 3$ from Born--Infeld Fibre Admissibility,''
    O23 spectral admissibility sub-programme paper, Zenodo (2026).
    \doi{10.5281/zenodo.19375136}.

    \bibitem{Q5a}
    J.~Beau,
    ``Large-$q$ Limit of the Admissible Fibre: Representational and Operator
    Convergence,''
    Cosmochrony companion paper~Q5a, Zenodo (2026).
    \doi{10.5281/zenodo.19642369}.

    \bibitem{Beau2026pto}
    J.~Beau,
    ``Projective Temporal Ordering and Cumulative Projected Capacity:
    Numerical Evidence from the Spectral Admissibility Cascade,''
    Cosmochrony spectral admissibility sub-programme, Zenodo (2026).
    \doi{10.5281/zenodo.19474639}.

    \bibitem{Bell1964}
    J.~S.~Bell,
    ``On the Einstein Podolsky Rosen paradox,''
    \textit{Physics Physique Fizika} \textbf{1}, 195--200 (1964).
    \doi{10.1103/PhysicsPhysiqueFizika.1.195}.

  \end{thebibliography}

\end{document}
